\section{Compact null families and descent}
\label{proj:sg:descent}

We return to the setting of Theorem~\ref{proj:thm:signed}, with
$0<v<n$, $r=v-1$, and the geometric construction of
Proposition~\ref{proj:prop:signed-construction}.  The infinitesimal
lifting result now turns a boundary fiber into a compact
subvariety outside the entire boundary.

\begin{proposition}\label{proj:prop:null-family}
Assume the transition surjectivity of
Proposition~\ref{proj:prop:formal-lifting}.  There is a dominating
algebraic family of integral projective subvarieties of $X$
of dimension
\[
 d=n-1-r=n-v
\]
whose general members avoid $D$ and on which $L$ is
numerically trivial.
\end{proposition}

\begin{proof}
Choose a separated affine etale chart $U\to\mathcal Y$
around a general point of a top-dimensional component of
$\mathcal T_U$.  Shrink it so that $C$ is trivial, this
component of $\mathcal T_U$ is smooth of pure dimension $r$,
and $S_U$ is flat over it.  We may also ensure that $S_U$
does not meet $J_U$, because the image of $J$ in $P$ has
dimension less than $r$.  Choose a closed point $t$ of this
open set and regular parameters $t_1,\ldots,t_r$ on
$\mathcal T_U$ cutting out $t$ alone after further shrinking.
Then
\[
 F=S_{U,t}
\]
is projective of pure dimension $d$, and the pulled-back
parameters form a regular sequence along $F$.

Put $I=\OO_{\mathcal X}(-S)$.  Proposition~\ref{proj:prop:formal-lifting}
gives surjections between all the sheaves
$g_*(I^j/I^{j+k})$ as the length $k$ increases.
Their transition kernels are coherent sheaves of the form
$g_*\OO_{S_U}\otimes C^{-j-k}$.  Since $U$ is affine, the
surjections hold on global sections as well.  We may
therefore lift the parameters compatibly to all
thickenings and lift the chosen conormal frame to a
compatible element $\widetilde y$ generating $I$ formally.

On $qS_U$, cut out the lifted $r$ parameters and denote the
resulting closed subscheme by $Z_q$.  Give it the structure
of a scheme over
\[
 R_q=\C[s]/(s^q),\qquad s\longmapsto\widetilde y.
\]
These schemes are compatible under reduction in $q$.
Their closed fiber is $F$, and they are flat over $R_q$.
Indeed the powers of the Cartier generator identify the
successive $s$-layers of $qS_U$ with $\OO_{S_U}$, proving
flatness before cutting.  The local flatness criterion then
shows that quotienting by lifts of a closed-fiber regular
sequence preserves flatness.

The map
$Z_q\to X\times\Spec R_q$ is quasi-finite.
It is proper as well: its reduction is the map from the
projective scheme $F$, and properness of a finite-type
morphism is unchanged by nilpotent thickenings.  Hence
this map is finite.  Let $R=\C[[s]]$.  Projective
Grothendieck existence, including its full faithfulness,
algebraizes the compatible finite algebra sheaves on
$X\times\Spec R_q$ to a finite algebra on
$X_R=X\times\Spec R$; multiplication and the unit
algebraize by full faithfulness
\cite[Lemmas~30.24.1 and~30.24.3]{Proj-StacksExistence}.
Its relative spectrum is a finite morphism
\[
 Z\longrightarrow X_R
\]
with the prescribed reductions.  The scheme $Z$ is
projective over $R$.  Formal flatness proves flatness
along the closed fiber, and flatness is automatic on the
generic fiber over $\C((s))$, so $Z$ is flat over $R$.

The images of the pole and coefficient-zero boundary
parts do not meet $F$.  Their inverse images in $Z$ are
closed and proper over $R$, so they do not meet $Z$:
a nonempty closed subset of a proper $R$-scheme has
specialization in the closed fiber.  Away from these parts
and the graph exceptional divisor, the root construction
identifies the divisor of $s_0$ with $\ell S$.
The formal generator $\widetilde y$ therefore supplies
compatible trivializations in which
\[
 N_0|_{\widehat Z}\simeq\OO_{\widehat Z},
 \qquad s_0|_{\widehat Z}=s^\ell.
\]
Full faithfulness algebraizes this trivialization and the
identity.  Thus the generic fiber of $Z$ avoids also the
positive part of $D$, and hence all of $D$.
Its finite image in $X_{\C((s))}$ has pure dimension $d$.
After a finite extension of $\C((s))$, a component can be
reduced and chosen geometrically integral.

It remains to show that these compact subvarieties cover
$X$, rather than merely a neighborhood of one point.
Choose a positive component $D_i$ with image dimension
$r$, and a component of $S$ covering it.  As the general
chart and the point $t$ vary, images of points of $F$
contain a dense open subset of $D_i$.
Every such point is in the specialization of a
$d$-dimensional generic finite image constructed above.
Indeed flatness over the discrete valuation ring rules
out vertical components of $Z$, and the dimension
formula shows that the closures of its generic
components have special-fiber components of dimension
$d$.  Their finite images retain that dimension.

Consider now all Hilbert schemes of $d$-dimensional
subvarieties of $X$.  Their geometrically integral loci
whose members avoid $D$ admit a countable stratification
by integral parameter spaces with irreducible universal
families.  Let $W_\alpha$ be the irreducible closures in
$X$ of the corresponding evaluation images.  A generic
image over a field extension determines a Hilbert point
of one of these strata.  Properness of the Hilbert scheme
gives its specialization, so every point of the dense
open subset of $D_i$ just described lies in some
$W_\alpha$.

Over the uncountable field $\C$, an irreducible variety
cannot have a dense open covered by countably many proper
closed subsets.  Therefore one $W_\alpha$ contains
$D_i$.  It also contains points outside $D$, by its
definition.  Since $D_i$ is a prime divisor in the
integral variety $X$, the only proper irreducible closed
subset containing it is $D_i$ itself.  Consequently
$W_\alpha=X$, giving a dominating family.
On every member of this family, the rational section
$s_0$ is regular and nowhere zero, since the member
avoids $D$.  It trivializes $N_0$, so $L$ is numerically
trivial there.
\end{proof}

\Needspace{12\baselineskip}
\subsection{Descent along the nef reduction}

\begin{lemma}\label{proj:sg:vertical-descent}
Let $f:Y\to B$ be a surjective projective morphism from a
normal integral variety to a smooth projective variety,
with geometrically connected integral generic fiber.
Suppose all fibers have dimension at most
$\dim Y-\dim B$.  Let $G_Y$ be a vertical $\Q$-Cartier
divisor that is relatively nef.  Then
\[
 G_Y=f^*G_B
\]
for a $\Q$-divisor $G_B$ on $B$.
\end{lemma}

\begin{proof}
If $B$ is a point, a vertical divisor is zero.  If the
relative dimension is zero, the fiber bound makes $f$
finite; its geometrically integral connected generic
fiber makes it birational, and normality of $B$ makes
it an isomorphism.  Both cases are immediate.  Assume
henceforth that the base and the relative dimension
are positive.

The fiber bound implies that every vertical prime divisor
maps onto a prime divisor of $B$.  For each such base
prime $P$, write the full pullback as
$f^*P=\sum_jm_jE_j$ and let $b_j$ be the coefficient of
$G_Y$ along $E_j$, including zero coefficients.
Give $P$ coefficient $\min_j(b_j/m_j)$ in $G_B$.
Only finitely many coefficients are nonzero.  Then
\[
 R'=G_Y-f^*G_B
\]
is effective, vertical, $\Q$-Cartier, relatively nef,
and misses at least one component over each base prime.
We prove that $R'=0$.

If $R'\ne0$, choose a base prime below its support.
Take a general complete-intersection curve in $B$
meeting that prime at a general point; when $\dim B=1$
use $B$ itself.  Its inverse image in $Y$ is normal
by normal Bertini.  It is integral: generic geometric
integrality gives the dominating component, and the
fiber bound excludes additional vertical components.
Next take $\dim Y-\dim B-1$ general very ample
hyperplanes upstairs.  This produces a normal integral
surface over the base curve with geometrically connected
generic fiber.  At the chosen general base point,
the cuts retain curves in the residual support and in
a missing component of the full pullback; these curves
are distinct because the original components were
distinct at the generic point of the base prime.

Resolve the surface.  The pulled-back residual divisor
is effective, vertical, and relatively nef.  Its
intersection with the full fiber is zero, whereas its
intersection with each fiber component is nonnegative.
Every one of the latter intersections is therefore zero.
The intersection matrix of a connected surface fiber
is negative semidefinite with kernel generated by the
full fiber.  Thus the residual part at the chosen point
is a multiple of that full fiber.  A missing component
forces this multiple to be zero, contradicting the
component in its support.  Connectedness of the fiber
follows from generic connectedness and Stein
factorization over the normal base curve.
\end{proof}

\begin{proof}[Completion of the proof of Theorem~\ref{proj:thm:signed}]
Lemma~\ref{proj:sg:numerical-boundary} handles $v=0$ and $v=n$.
In the remaining case, combine
Propositions~\ref{proj:prop:signed-construction},
\ref{proj:prop:formal-lifting}, and~\ref{proj:prop:null-family}.

Apply the nef-reduction theorem to a Cartier multiple
of $L$ \cite[Theorem~2.1]{Proj-BauerEtAl2002}.  It gives an almost
holomorphic rational map with connected fibers, with
$L$ numerically trivial on its compact general fibers,
and with positive $L$-degree on every noncontracted
curve through a very general point of $X$.
Let its base dimension be $b$.  A compact $L$-trivial
$d$-fold through such a point is contracted: otherwise
general ample slices through the point yield a
noncontracted curve of $L$-degree zero.  Hence
\[
 b\le n-d=v.
\]

Resolve the graph, flatten the main component over a
modification of the base, resolve that base, and
normalize the main transform.  We obtain projective
morphisms
\[
 \begin{tikzcd}
 Y \arrow[r,"\pi"] \arrow[d,"f"'] & X\\
 B &
 \end{tikzcd}
\]
with $\pi$ birational, $B$ smooth, $Y$ normal, and
geometrically connected integral generic fiber of $f$.
All fibers have dimension at most $n-b$.
Indeed the flat main transform remains integral after
the base modification by flatness and generic
integrality, and finite normalization preserves the
fiber-dimension bound.  Normalization need not preserve
flatness; only this bound is used.

The pullback $\pi^*D$ has no horizontal component.
To see this, restrict to a very general fiber of $f$.
The class of $\pi^*L$ is numerically trivial there,
and $\pi^*(L-cD)$ restricts to a pseudo-effective class.
The latter assertion follows by restricting effective
approximants, avoiding the countably many fibers
contained in their supports.  Its restricted class is
$-c\pi^*D$.  A negative nonzero effective divisor on a
projective variety cannot be pseudo-effective, as
intersection with an ample power shows.  Thus the
restriction of $\pi^*D$ is zero.  In particular $b=0$
would force $D=0$, contrary to $v>0$.

Consequently $G_Y=\pi^*G$ is vertical and is relatively
nef, since $G_Y\sim_{\Q}\pi^*L$.  Lemma~\ref{proj:sg:vertical-descent}
gives
\[
 \pi^*L\sim_{\Q}f^*G_B
\]
for a $\Q$-divisor $G_B$ on the smooth base.
It is nef: every curve of $B$ is dominated by a curve
of $Y$, on which the pullback has nonnegative degree.
The intersection formula for a pullback gives
$\nu(X,L)\le b$.  Combined with $b\le v$, this yields
$b=v$ and $G_B^b>0$.  Therefore $G_B$ is big.
Pulling back sections gives
$\kappa(X,L)\ge b=v$, and the general inequality
$\kappa(X,L)\le\nu(X,L)$ for a nef divisor proves
abundance.
\end{proof}
